% !PRINTED 210
\weberpage{210}
%<<CONT
und daraus, durch Auflösung, mit Benutzung der Bezeichnung \S\, 28, (4):

\weq{5}{m(\mu, \mu') \equiv \begin{pmatrix} a, & b \\ c, & \partial \end{pmatrix} (\nu, \nu').}

Es ist daher

\weq{6}{\begin{pmatrix} a, & b \\ c, & \partial \end{pmatrix}}

ein zweckmäßiges Zeichen für die Vertauschung $S$, und aus (5) ersieht man, daß sich zwei solche Vertauschungen:

\[
S = \begin{pmatrix} a, & b \\ c, & \partial \end{pmatrix}, \qquad S' = \begin{pmatrix} a', & b' \\ c', & \partial' \end{pmatrix}
\]

genau nach der in \S\, 28 gegebenen Regel:

\weq{7}{S S' = S'' = \begin{pmatrix} a a' + b c', & a b' + b \partial' \\ c a' + \partial c', & c b' + \partial \partial' \end{pmatrix}}

zusammensetzen, so daß die Vertauschung $S''$ der Wurzeln entsteht, wenn \gesperrt{zuerst} die Vertauschung $S$, \gesperrt{sodann} die Vertauschung $S'$ unter den Wurzeln der Teilungsgleichung vorgenommen wird. Zu beachten ist aber immer, daß hier nur die nach dem Modul $n$ genommenen Reste der Zahlen $a, b, c, \partial$ in Betracht kommen, so daß die Anzahl der Vertauschungen $S$ stets endlich ist\footnote{Wollte man, was auf den ersten Blick näher zu liegen scheint, die Bezeichnung $S = \begin{pmatrix} \partial, & -b \\ -c, & a \end{pmatrix}$ wählen, so würde die Komposition nach der Formel $S'S = S''$, also umgekehrt wie bei der üblichen Komposition der Permentationen % [sic]
, zu bezeichnen sein.}.

1. \gesperrt{Der Inbegriff aller Substitutionen $\begin{pmatrix} a, & b \\ c, & \partial \end{pmatrix}$ bildet eine Gruppe, die wir mit $\mathfrak{A}$ bezeichnen, und in ihr ist, wie wir bewiesen haben, die Gruppe $\mathfrak{G}$ der Teilungsgleichung enthalten.}

In $\mathfrak{A}$ ist als Teiler eine Gruppe $\mathfrak{B}$ enthalten, die aus allen Substitutionen (5) % [sic] The source's cross-reference appears to require (8) or (6).

\weq{8}{T = \begin{pmatrix} \alpha, & \beta \\ \gamma, & \delta \end{pmatrix}}

besteht, die der Bedingung

\weq{9}{\alpha\delta - \beta\gamma \equiv 1 \pmod{n}}

genügen.
